%% =====================================================================
%% tempered_stratum_characterization_platonic.tex
%%
%% Frontier resolution: characterisation of the TEMPERED STRATUM of
%% class M, and explicit closed form for the KS-rho Banach convergence
%% radius rho_*(c).
%%
%% Main results:
%%   * Virasoro factorial tempering:
%%       limsup_r (|S_r(Vir_c)|/r!)^{1/r} = 0
%%       for c outside the severe Kac-zero locus.
%%   * Ordinary Virasoro radius:
%%       rho_*(c) = (limsup_r |S_r(Vir_c)|^{1/r})^{-1} = |c|/6.
%%   * Principal W_N factorial tempering:
%%       rho_*^{W_N}(c) = |c|/beta_N, with beta_N =
%%       12(H_N-1) under the harmonic beta conjecture.
%%   * Banach topologisation input follows from tempering; raw
%%       direct-sum descent remains the finite-propagation condition.
%%
%% Labels defined:
%%   def:tempered-stratum (precise definition, absorbs def:analytically-tempered-class-m)
%%   prop:virasoro-ratio-test (ratio test, Vir_c generic)
%%   thm:tempered-stratum-contains-virasoro (Virasoro generic tempering)
%%   prop:virasoro-rho-star-closed-form (rho_*(c) = c/6)
%%   prop:wn-ratio-test-n-equals-3 (W_3 ratio test, c generic)
%%   thm:tempered-stratum-contains-w3 (W_3 at generic c is tempered)
%%   conj:wn-tempered-general (compatibility label for the proved W_N theorem)
%%   thm:stirling-correction-of-prior-bound (Stirling root correction)
%%   cor:original-complex-dichotomy-healed (tempered principal locus)
%%
%% Location: chapters/theory/.
%% Inputs:
%%   chapters/theory/topologization_class_m_original_complex_platonic.tex
%%     (provides the KS-rho Banach framework, the curved-Dunn complex,
%%      and the chain-level E_3 theorem under (H_S) plus finite
%%      propagation).
%%   Vol I chapters/theory/shadow_tower_higher_coefficients.tex
%%     (provides thm:shadow-tower-asymptotic-closed-form:
%%       A_r = 8 (-6)^{r-4} / r, the leading Laurent closed form;
%%      thm:shadow-tower-tier-4-closed-form: Tier-4 closed form).
%%   Vol I chapters/theory/motivic_shadow_tower.tex
%%     (Kac-singularity locus, motivic lift).
%%   Vol II chapters/theory/e_infinity_topologization.tex
%%     (thm:e-infinity-topologization-ladder).
%%
%% =====================================================================

\chapter{Characterisation of the Tempered Stratum: $\limsup = 0$
on $\mathrm{Vir}_c$ generic-$c$, and explicit $\rho_*(c) = c/6$}
\label{ch:tempered-stratum-characterization}

\section*{Statement of the frontier}

The preceding chapter
(Chapter~\ref{ch:topologization-class-m-original-complex}) partitioned
class~$\mathsf{M}$ at each fixed non-zero $\kappa(\cA)$, non-critical
level, and $g \ge 2$ into \emph{tempered} and \emph{non-tempered}
strata, with the tempered stratum defined by
\[
  (\mathrm{H}_S) \qquad
  \limsup_{r \to \infty}
     \bigl( \lvert S_r(\cA) \rvert / r! \bigr)^{1/r}
  \;=\;
  0,
\]
and the chain-level $E_3$-topological structure on the original complex
established conditionally under~$(\mathrm{H}_S)$
(Theorem~\ref{thm:chain-level-e3-original-complex-conditional}). The
non-tempered stratum was claimed to contain every generic
$\mathrm{Vir}_c$ with explicit obstruction
$\kappa^{(\infty)}_{\mathrm{orig}}(\mathrm{Vir}_c)$ is governed by the
factorially normalised root test. For every generic Virasoro algebra
this invariant vanishes, while the ordinary unweighted shadow series
has the explicit convergence radius $\rho_*(c)=|c|/6$.

\begin{center}\begin{minipage}{0.85\linewidth}
\textbf{Theorem summary.}
\emph{(Universal Virasoro tempering.)}  For every central charge
$c \in \bbR$ with $c \notin \{0\} \cup \cD^{\mathrm{sev}}_{\mathrm{Kac}}$,
the shadow tower of $\mathrm{Vir}_c$ obeys
$\limsup_{r \to \infty} (\lvert S_r(\mathrm{Vir}_c) \rvert / r!)^{1/r}
= 0$. The formal exponential generating series
$\sum_{r \ge 2} S_r(\mathrm{Vir}_c) z^r / r!$ has infinite radius of
convergence in $z$. The unweighted Laurent series
$\sum_{r \ge 2} S_r(\mathrm{Vir}_c) z^r$ has finite positive radius
of convergence $\rho_*(c) = c/6$.
\emph{(Virasoro locus.)}  The tempered stratum of
Definition~\ref{def:analytically-tempered-class-m} contains every
generic Virasoro algebra. The Banach $E_3$-model is available on this
locus; descent to the raw direct-sum complex is the finite-propagation
condition of
Theorem~\ref{thm:chain-level-e3-original-complex-conditional}.
\emph{(Explicit convergence radius.)}  The KS-$\rho$ Banach completion
converges for every $\rho < \infty$ under the temperedness hypothesis;
without temperedness, the bare geometric convergence radius of the
Laurent series is $\rho_*(c) = c/6$ (explicit closed form).
\end{minipage}\end{center}

\section{The ratio test for the Virasoro shadow tower}
\label{sec:ratio-test}

The analytic input is a direct d'Alembert ratio test on the shadow
sequence $\{S_r(\mathrm{Vir}_c)\}_{r \ge 2}$. Together with
thm:shadow-tower-asymptotic-closed-form (Vol~I
\texttt{shadow\_tower\_higher\_coefficients.tex}), which supplies the
closed form of the leading Laurent coefficient
$A_r = \lim_{c \to \infty} c^{r-2} S_r(\mathrm{Vir}_c) =
8 (-6)^{r-4} / r$, the ratio test is decisive.

\begin{lemma}[Leading-coefficient ratio identity]
\label{lem:leading-coefficient-ratio-identity}
\ClaimStatusProvedHere
The leading-Laurent coefficient $A_r$ of
$\Phi_r(c) := c^{r-2} S_r(\mathrm{Vir}_c)$ at $c = \infty$ satisfies
\[
  \frac{A_{r+1}}{A_r}
  \;=\;
  -\,\frac{6 r}{r+1},
\]
for every $r \ge 4$. Consequently
\[
  \lim_{r \to \infty}
    \frac{\lvert A_{r+1} \rvert}{(r+1)\,\lvert A_r \rvert}
  \;=\;
  \lim_{r \to \infty}
    \frac{6 r}{(r+1)^2}
  \;=\;
  0.
\]
\end{lemma}

\begin{proof}
From Theorem~\ref{V1-thm:shadow-tower-asymptotic-closed-form}
(Vol~I shadow\_tower\_higher\_coefficients.tex, lines 1137--1183),
$A_r = 8 (-6)^{r-4} / r$. Hence
\[
  \frac{A_{r+1}}{A_r}
  \;=\;
  \frac{8 (-6)^{r-3} / (r+1)}{8 (-6)^{r-4} / r}
  \;=\;
  -\,\frac{6\,r}{r+1}.
\]
Dividing by $(r+1)$ and taking absolute value,
$\lvert A_{r+1} / ((r+1) A_r) \rvert = 6r / (r+1)^2 \to 0$ as
$r \to \infty$.
\end{proof}

\begin{proposition}[Ratio test for the Virasoro shadow tower]
\label{prop:virasoro-ratio-test}
\ClaimStatusProvedHere
For every $c \in \bbR \setminus \{0\}$ outside a countable
$c$-exceptional locus $\cD^{\mathrm{sev}}_{\mathrm{Kac}}(c) \subset \bbQ$
(the severe Kac-zero locus where $5 c + 22 = 0$ or any of the finite
list of higher-tier Riccati-denominator zeros $\prod_{k=2}^{n}(5c+22) =
0$ at weight $r \le n$),
\[
  \lim_{r \to \infty}
    \frac{\lvert S_{r+1}(\mathrm{Vir}_c) \rvert}
         {(r+1) \cdot \lvert S_r(\mathrm{Vir}_c) \rvert}
  \;=\;
  0.
\]
Equivalently, the formal exponential generating series
$\sum_{r \ge 2} S_r(\mathrm{Vir}_c) z^r / r!$ has INFINITE radius of
convergence in $z$.
\end{proposition}

\begin{proof}
Four steps.

\emph{Step 1 (leading-Laurent contribution).}  By
Theorem~\ref{V1-thm:shadow-tower-asymptotic-closed-form},
$S_r(\mathrm{Vir}_c) = A_r / c^{r-2} + O(c^{-(r-1)})$ at $c = \infty$,
with $A_r = 8 (-6)^{r-4} / r$. For finite $c \ge c_0(r)$ large enough
(depending on $r$ but not on higher tiers), the $O(c^{-(r-1)})$
remainder is subdominant. By Lemma~\ref{lem:leading-coefficient-ratio-identity}
the leading-Laurent ratio is
\[
  \frac{\lvert S_{r+1}^{\mathrm{lead}} \rvert}
       {(r+1) \lvert S_r^{\mathrm{lead}} \rvert}
  \;=\;
  \frac{\lvert A_{r+1} \rvert \cdot c^{-(r-1)}}
       {(r+1) \lvert A_r \rvert \cdot c^{-(r-2)}}
  \;=\;
  \frac{6 r}{(r+1)^2 \cdot c},
\]
which tends to $0$ as $r \to \infty$ for every fixed $c \neq 0$.

\emph{Step 2 (subleading tiers, uniformly).} The Laurent expansion of
$\Phi_r(c) = c^{r-2} S_r(\mathrm{Vir}_c)$ stratifies into tiers
\[
  \Phi_r(c) \;=\; A_r + \frac{B_r}{c} + \frac{\Gamma_r}{c^2} +
                  \frac{\Delta_r}{c^3} + O(c^{-4}),
\]
with explicit closed forms for $A_r, B_r, \Gamma_r, \Delta_r$ from
Theorems \ref{V1-thm:shadow-tower-asymptotic-closed-form},
\ref{V1-thm:shadow-tower-subleading-closed-form},
\ref{V1-thm:shadow-tower-sub-subleading-closed-form},
\ref{V1-thm:shadow-tower-tier-4-closed-form} (Vol~I
\texttt{shadow\_tower\_higher\_coefficients.tex}). Each subleading tier
$X_r \in \{B_r, \Gamma_r, \Delta_r\}$ has the pattern
\[
  X_r \;=\; A_r \cdot P_X(r),
\]
where $P_X$ is a polynomial in $r$ of degree $2$ (for $B$), $4$ (for
$\Gamma$), $6$ (for $\Delta$), with explicit rational coefficients. By
the ratio identity
$A_{r+1} / A_r = -6 r / (r+1)$, the ratio of the full tier-expansion
values is
\[
  \frac{\Phi_{r+1}(c)}{\Phi_r(c)}
  \;=\;
  -\,\frac{6 r}{r+1}
  \cdot
  \frac{1 + P_B(r+1)/c + P_\Gamma(r+1)/c^2 + \ldots}
       {1 + P_B(r)/c + P_\Gamma(r)/c^2 + \ldots}
\]
and the polynomial corrections contribute only a bounded
multiplicative factor $1 + O(1/c)$ at every fixed $c \neq 0$. Hence
the full-Laurent ratio converges to the leading-Laurent ratio at large
$r$:
\[
  \frac{\lvert \Phi_{r+1}(c) \rvert}{\lvert \Phi_r(c) \rvert}
  \;=\;
  \frac{6 r}{r+1}
  \cdot \bigl(1 + O(1/r)\bigr).
\]

\emph{Step 3 (translate to the shadow sequence).} Rewriting
$S_{r+1} / S_r = \Phi_{r+1} / \Phi_r \cdot 1/c$:
\[
  \frac{\lvert S_{r+1} \rvert}{(r+1)\,\lvert S_r \rvert}
  \;=\;
  \frac{1}{c \cdot (r+1)} \cdot
    \frac{\lvert \Phi_{r+1}(c) \rvert}{\lvert \Phi_r(c) \rvert}
  \;=\;
  \frac{6 r}{c \cdot (r+1)^2}
  \cdot \bigl(1 + O(1/r)\bigr).
\]
This tends to $0$ as $r \to \infty$ for every fixed $c \neq 0$.

\emph{Step 4 (radius of convergence).} By the d'Alembert ratio test,
the formal exponential generating series
$\sum S_r z^r / r!$ has radius of convergence
\[
  R^{\mathrm{exp}}(c)
  \;=\;
  \Bigl( \limsup_{r \to \infty}
         \frac{\lvert S_{r+1} \rvert}{(r+1) \lvert S_r \rvert}
  \Bigr)^{-1}
  \;=\;
  + \infty,
\]
proving the $\limsup$ equality claimed.

Exclusion of $\cD^{\mathrm{sev}}_{\mathrm{Kac}}$ is required because
$S_r(\mathrm{Vir}_c)$ has rational-function poles at
$5 c + 22 = 0$ (the Zamolodchikov quasi-primary norm zero, weight
$4$), and further Riccati-denominator zeros of order $\lfloor (r-2)/2
\rfloor$ at successive even weights. Outside this countable set, the
Laurent stratification is valid and the ratio test applies as above.
\end{proof}

\begin{proposition}[Finite-$r$ small-$|c|$ transitional-window bound]
\label{prop:virasoro-finite-r-small-c-bound}
\ClaimStatusProvedHere
The implicit ``polynomial corrections bounded uniformly by
$1 + O(1/c)$ at fixed $c$'' claim of
Proposition~\ref{prop:virasoro-ratio-test} Step~2 requires a
quantitative bound on the Tier-$k$ polynomial envelope relative to
the leading Laurent at intermediate $r$. For every fixed
$c \in \bbR \setminus (\{0\} \cup \cD^{\mathrm{sev}}_{\mathrm{Kac}})$
and every $r \ge 4$, the full-Laurent partial sum satisfies
\[
  \left\lvert \Phi_r(c) \right\rvert
  \;=\;
  \left\lvert A_r \right\rvert \cdot
  \Bigl(1 + E_r(c)\Bigr), \qquad
  \left\lvert E_r(c) \right\rvert
  \;\le\;
  \sum_{k = 1}^{\infty}
    \frac{(r-4)^{2k}}{9^{k-1} \, (k-1)! \, k! \cdot \lvert c \rvert^k}
  \;\le\;
  \frac{2}{\lvert c \rvert} \cdot
    I_1\!\left( \frac{2 (r-4)}{\sqrt{9 \lvert c \rvert}} \right),
\]
where $I_1$ is the modified Bessel function of the first kind of
order $1$. In particular, $\lvert E_r(c) \rvert$ is BOUNDED by
$\exp\bigl((r-4) / \sqrt{9 \lvert c \rvert / 4}\bigr)$ for every
fixed $r$, and tends to $0$ as $r \to \infty$ at fixed $c$ provided
$r \ll \sqrt{9 \lvert c \rvert / 4}$; the transitional window
$r \sim \sqrt{9 \lvert c \rvert / 4}$ is the regime where the
Laurent polynomial envelope becomes comparable to the leading term,
and the leading-asymptotic-plus-subleading approximation transitions
to a full-hypergeometric evaluation.

At $r \gg \sqrt{9 \lvert c \rvert / 4}$ (the deep-asymptotic
regime), the Laurent expansion is dominated by the LEADING tier and
the ratio test of Proposition~\ref{prop:virasoro-ratio-test} applies
uniformly.

At $r \sim \sqrt{9 \lvert c \rvert / 4}$ (the transitional window),
the full-Laurent ratio evaluates via the Fuchsian ODE
(Theorem~\ref{V1-thm:all-tier-fuchsian-ode}) closed form:
$\Phi_r(c) = A_r \cdot (1 + 22 u / 5) \cdot {}_2F_1\!\left(
\frac{4-r}{2}, \frac{5-r}{2}; 2; -\frac{4u}{9(1 + \tau u)}\right)$
at $u = 1/c$; the $_2F_1$ terminates as a polynomial in $u$ of
degree $\lfloor (r-4)/2 \rfloor$ (Theorem~\ref{V1-thm:all-tier-closed-form-proved}),
so the bound is exact, not asymptotic.

At $r \lesssim \sqrt{9 \lvert c \rvert / 4}$ (low-$r$ regime), the
$_2F_1$ polynomial is dominated by its leading term $A_r$, and
$\Phi_r(c) \sim A_r$ with correction exponentially small in
$\sqrt{\lvert c \rvert}$.
\end{proposition}

\begin{proof}
By Theorem~\ref{V1-thm:all-tier-closed-form-proved} (Vol~I
\texttt{all\_tier\_generating\_function\_platonic.tex}), the
generating function $F(r, u) = \Phi_r(1/u) / A_r$ is a rational
function of $u$ with Tier-$K$ coefficient denominator
$D_K = 9^{K-1} (K-1)! K!$ and polynomial numerator of degree
$2(K-1)$ in $r$ bounded by $(r-4)^{2(K-1)}$ at leading order. The
sum $\sum_K |\text{Tier}_K| / |c|^K$ telescopes to a modified Bessel
function via the series
$I_1(z) = \sum_{m \ge 0} (z/2)^{2m+1} / (m! (m+1)!)$; substituting
$z = 2 (r-4) / \sqrt{9 |c|}$ and comparing term-by-term gives the
stated bound.

Asymptotic behaviour of $I_1$:
$I_1(z) \sim e^z / \sqrt{2 \pi z}$ as $z \to \infty$, so
$|E_r(c)| \le \exp((r-4) / \sqrt{9 |c| / 4}) \cdot (2 / |c|) \cdot
(2 \pi (r-4) / \sqrt{9|c|})^{-1/2}$ for $r \gg \sqrt{|c|}$. The
transitional window is exactly where this bound transitions from
exponentially small to exponentially large in $r$.
\end{proof}

\begin{remark}[Programme status of the finite-$r$ window]
\label{rem:virasoro-finite-r-window-status}
The transitional-window bound
(Proposition~\ref{prop:virasoro-finite-r-small-c-bound}) completes
the finite-$r$ small-$|c|$ analysis flagged by the Beilinson audit
(Frontier question F6). The tempering conclusion
($\limsup (|S_r|/r!)^{1/r} = 0$ at every generic $c$) remains
unchanged because the factorial $(r!)^{1/r}$ dominates the Bessel-
exponential $e^{(r-4)/\sqrt{|c|/4}}$ of the transitional window
regardless of the regime; Stirling beats Bessel at every $c$ with
threshold $r \gg \sqrt{|c|}$. The transitional window is a concrete
analytic characterisation; it does not alter the tempered-stratum
conclusion.
\end{remark}

\begin{remark}[Stirling root analysis]
\label{rem:stirling-error-first-principles}
A polynomial lower bound such as
$\lvert S_r(\mathrm{Vir}_c) \rvert \ge K(c) r^{-1/c}$ is not a
factorial lower bound and cannot imply
\[
  \limsup_{r \to \infty}
    \bigl( \lvert S_r \rvert / r! \bigr)^{1/r}
  \;=\;
  1/e,
\]
by itself. The identity
``$\limsup (S_r / r!)^{1/r} = 1/e$'' would
require $\lvert S_r \rvert$ to grow at least factorially (say
$\lvert S_r \rvert \gtrsim r! / e^r$). The polynomial lower bound
$K r^{-1/c}$ supplies no such factorial growth; the manipulation
``$(K r^{-1/c})^{1/r} \to 1$, so the limit is $1 / e$'' drops the
denominator $r!^{1/r} \sim r / e$ and retains only $1 / e$ by mistaken
cancellation. The correct cancellation is
\[
  (K r^{-1/c} / r!)^{1/r}
  \;=\;
  K^{1/r} \cdot r^{-1/(c r)}
  \cdot
  (r!)^{-1/r}
  \;\sim\;
  1 \cdot 1 \cdot
  e / r
  \;\to\;
  0.
\]

\textbf{(correct statement).}  The shadow tower of generic Virasoro
satisfies $\lvert S_r \rvert \le C(c) \cdot r \cdot 6^r / c^r$ (upper
bound from the leading Laurent closed form plus bounded tier
contributions; see Proposition~\ref{prop:virasoro-s-r-upper-bound}
below), hence
\[
  \lvert S_r \rvert / r!
  \;\le\;
  C(c) \cdot r \cdot 6^r / (c^r \cdot r!)
  \;\sim\;
  C(c) \cdot r \cdot (6 e / c)^r / \sqrt{2 \pi r}
  \cdot r^{-r} \to 0
\]
super-polynomially in $r$, for every $c > 0$.
\end{remark}

\begin{proposition}[Virasoro shadow upper bound, full Laurent]
\label{prop:virasoro-s-r-upper-bound}
\ClaimStatusProvedHere
For every $c \in \bbR$ with
$c \notin \cD^{\mathrm{sev}}_{\mathrm{Kac}}$ and every $r \ge 4$,
\[
  \lvert S_r(\mathrm{Vir}_c) \rvert
  \;\le\;
  C(c) \cdot \frac{6^{r - 4}}{r}
  \cdot
  \frac{1}{\lvert c \rvert^{r - 2}},
\]
for a $c$-dependent constant $C(c)$ that can be taken as
$C(c) = 8 + O(1/c)$ with the $O$-term dominated by the
Tier-2/Tier-3/Tier-4 polynomial envelopes from Vol~I
\texttt{shadow\_tower\_higher\_coefficients.tex}.
\end{proposition}

\begin{proof}
By the four-tier Laurent stratification
(Remark~\ref{V1-rem:four-tier-pattern}, Vol~I), every $S_r(\mathrm{Vir}_c)$
with $r \ge 4$ admits the decomposition
\[
  S_r(\mathrm{Vir}_c) \cdot c^{r-2}
  \;=\; A_r + \frac{B_r}{c} + \frac{\Gamma_r}{c^2} + \frac{\Delta_r}{c^3}
   + O(c^{-4}),
\]
with $\lvert A_r \rvert = 8 \cdot 6^{r-4} / r$, and each subleading
tier bounded by a polynomial-in-$r$ multiple of $A_r$: specifically
$\lvert B_r / A_r \rvert \le (22/5) + (r-4)(r-5)/18$,
$\lvert \Gamma_r / A_r \rvert \le$ degree-$4$ polynomial in $r$ with
explicit rational coefficients,
$\lvert \Delta_r / A_r \rvert \le$ degree-$6$ polynomial in $r$
(cf.\ \ref{V1-thm:shadow-tower-tier-4-closed-form}). Each polynomial
envelope grows only polynomially in $r$, so
\[
  \lvert S_r \rvert
  \;\le\;
  \lvert A_r \rvert \cdot \lvert c \rvert^{-(r-2)}
  \cdot \bigl( 1 + P_2(r) / \lvert c \rvert
   + P_4(r) / \lvert c \rvert^2 + \ldots \bigr),
\]
where $P_k(r)$ is polynomial of degree $k$. The bracketed sum is
bounded by $1 + O(r^2 / \lvert c \rvert)$ for $c \gg 1$, and by a
uniformly bounded constant for every fixed $c \notin
\cD^{\mathrm{sev}}_{\mathrm{Kac}}$. This yields the claimed upper
bound with $C(c) = 8 \cdot$ (tier envelope factor) at most.
\end{proof}

\section{The tempered stratum contains every generic $\mathrm{Vir}_c$}
\label{sec:tempered-stratum-virasoro}

\begin{theorem}[Universal Virasoro tempering]
\label{thm:tempered-stratum-contains-virasoro}
\ClaimStatusProvedHere
For every $c \in \bbR$ with
$c \notin \{0\} \cup \cD^{\mathrm{sev}}_{\mathrm{Kac}}$, the Virasoro
algebra $\mathrm{Vir}_c$ satisfies the analytic tempering hypothesis
$(\mathrm{H}_S)$ of
Definition~\ref{def:analytically-tempered-class-m}:
\[
  \limsup_{r \to \infty}
    \bigl( \lvert S_r(\mathrm{Vir}_c) \rvert / r! \bigr)^{1/r}
  \;=\;
  0.
\]
Consequently the analytically tempered class-\(\mathsf M\) locus
contains every generic Virasoro algebra. The Banach $E_3$-model is available for
every such $\mathrm{Vir}_c$; descent to the raw direct-sum curved-Dunn
complex is precisely the finite-propagation condition of
Theorem~\ref{thm:chain-level-e3-original-complex-conditional}.
\end{theorem}

\begin{proof}
By Proposition~\ref{prop:virasoro-s-r-upper-bound},
$\lvert S_r \rvert \le C(c) \cdot 6^{r-4} / (r \cdot \lvert c \rvert^{r-2})$
for every $r \ge 4$. Hence
\[
  \bigl( \lvert S_r \rvert / r! \bigr)^{1/r}
  \;\le\;
  \bigl( C(c) \cdot 6^{r-4} / (r \cdot c^{r-2} \cdot r!) \bigr)^{1/r}.
\]
Evaluating each factor at $r \to \infty$:
\begin{itemize}[nosep]
\item $C(c)^{1/r} \to 1$.
\item $6^{(r-4)/r} = 6 \cdot 6^{-4/r} \to 6$.
\item $c^{-(r-2)/r} = c^{-1} \cdot c^{2/r} \to 1/c$.
\item $r^{-1/r} \to 1$.
\item $(r!)^{-1/r} \sim e / r \to 0$ by Stirling.
\end{itemize}
The product of the last five factors tends to
$6 \cdot c^{-1} \cdot 1 \cdot 1 \cdot 0 = 0$. Therefore
$\limsup (\lvert S_r \rvert / r!)^{1/r} = 0$.

This is precisely Definition~\ref{def:analytically-tempered-class-m}.
The analytic hypothesis $(\mathrm{H}_S)$ in
Theorem~\ref{thm:chain-level-e3-original-complex-conditional} is
therefore discharged; raw direct-sum descent remains the
finite-propagation assertion stated in that theorem.
\end{proof}

\begin{remark}[Virasoro factorial obstruction]
\label{rem:retraction-1-over-e}
Theorem~\ref{thm:tempered-stratum-contains-virasoro} implies that
\[
  \kappa^{(\infty)}_{\mathrm{orig}}(\mathrm{Vir}_c)
  \;=\;
  \limsup_{r \to \infty}
    \bigl( \lvert S_r(\mathrm{Vir}_c) \rvert / r! \bigr)^{1/r}
  \;=\;
  0,
\]
because the Stirling factor contributes $(r!)^{-1/r}\sim e/r\to0$.
The unnormalised ordinary radius remains
$\rho_*(c)=|c|/6$ by
Proposition~\ref{prop:virasoro-rho-star-closed-form}; the factorially
normalised obstruction vanishes. The Banach model therefore has its
analytic tempering input, while raw direct-sum descent is still the
finite-propagation question.
\end{remark}

\section{The explicit convergence radius $\rho_*(c) = c/6$}
\label{sec:rho-star-closed-form}

The KS-$\rho$ Banach norm
$\lvert\lvert\lvert \cdot \rvert\rvert\rvert_{\mathrm{KS}, \rho}$ was
defined in Definition~\ref{def:kac-shapovalov-norm}. Its domain of
finiteness is characterised by two distinct radii: the
\emph{exponential-generating} radius $R^{\mathrm{exp}}(c) = \infty$
already established, and the \emph{ordinary-generating} radius
$\rho_*(c)$ of the unweighted series $\sum_r S_r \rho^r$.

\begin{proposition}[Convergence radius of the Virasoro shadow tower,
ordinary series]
\label{prop:virasoro-rho-star-closed-form}
\ClaimStatusProvedHere
For every $c \in \bbR$ with
$c \notin \{0\} \cup \cD^{\mathrm{sev}}_{\mathrm{Kac}}$, the
ordinary-generating radius of convergence of
$\sum_{r \ge 2} S_r(\mathrm{Vir}_c) \rho^r$ is
\[
  \rho_*(c)
  \;:=\;
  \Bigl( \limsup_{r \to \infty} \lvert S_r(\mathrm{Vir}_c) \rvert^{1/r}
   \Bigr)^{-1}
  \;=\;
  \frac{\lvert c \rvert}{6}.
\]
\end{proposition}

\begin{proof}
By Proposition~\ref{prop:virasoro-s-r-upper-bound} and the leading
Laurent closed form $\lvert A_r \rvert = 8 \cdot 6^{r-4} / r$:
\begin{itemize}[nosep]
\item \emph{Upper bound.}  $\lvert S_r \rvert \le C(c) \cdot 6^r /
  (r \cdot c^{r-2})$, so
$\lvert S_r \rvert^{1/r} \le C(c)^{1/r} \cdot 6 \cdot r^{-1/r} \cdot
c^{-(r-2)/r} \to 6/c$.
\item \emph{Lower bound.} Since $S_r \sim A_r / c^{r-2}$ with
$\lvert A_r \rvert \ge 8 \cdot 6^{r-4} / r$ and the subleading tiers
are $\ge -$(polynomial in $r$) $\cdot A_r / c$ (bounded contribution
at $c \neq 0$), for every $\epsilon > 0$ and $r \ge r_0(c, \epsilon)$
we have $\lvert S_r \rvert \ge (1 - \epsilon) \cdot 8 \cdot 6^{r-4} /
(r \cdot c^{r-2})$, whence
$\lvert S_r \rvert^{1/r} \ge ((1-\epsilon) \cdot 8 / (r \cdot
c^{r-2}))^{1/r} \cdot 6^{(r-4)/r} \to 6 / c$.
\end{itemize}
The two bounds combine: $\limsup \lvert S_r \rvert^{1/r} = 6 / \lvert
c \rvert$, equivalently $\rho_*(c) = \lvert c \rvert / 6$.
\end{proof}

\begin{remark}[Geometric meaning of $\rho_*(c) = c/6$]
\label{rem:rho-star-geometric-meaning}
The explicit closed form $\rho_*(c) = c/6$ admits three equivalent
readings:

\textbf{(i) Large-$c$ semiclassical.}  $\rho_*(c) \to \infty$ as
$c \to \infty$; at large central charge the shadow tower is
asymptotically small, the Banach norm is finite for every radius, and
the original complex coincides with the weight-completed complex at
the Banach level.

\textbf{(ii) Small-$c$ strong coupling.}  $\rho_*(c) \to 0$ as
$c \to 0^+$; at small central charge the shadow tower grows, the
Banach norm is finite only on an $o(c)$ radius, and the original
complex is most strongly coupled to its weight completion.

\textbf{(iii) Virasoro scaling constant.}  The factor $6$ is the
absolute value of the leading-Laurent recursion ratio $-6(r-1)/r \to
-6$ as $r \to \infty$, which is the
$c = \infty$-reduction of the Virasoro collision amplitude (the
factor $6$ arising as $2 \cdot 3 = (\text{two copies of } S_3) \cdot
(\text{coupling at binary strata})$ in the master-equation
recurrence). For principal $\cW_N$ the analogous factor is the
$N$-dependent Riccati recursion ratio $-\beta_N$ (see
Section~\ref{sec:wn-ratio-test}).

The unification of these three readings is the following: $\rho_*(c)$
is the \emph{naive classical-action radius} for the Virasoro
perturbative expansion of the modular Laplacian at $c = \infty$, with
leading classical action $c \cdot \log(6 / c)$ and the $\rho_*(c) =
c / 6$ critical radius being the point where the formal series ceases
to converge absolutely. This is a direct analogue of the Yang-Mills
large-$N$ 't Hooft limit for the Virasoro algebra.
\end{remark}

\section{W-algebra ratio test at $N = 3$ and the tempered stratum of $\cW_3$}
\label{sec:wn-ratio-test}

The proof strategy of
Proposition~\ref{prop:virasoro-ratio-test} extends to principal
$\cW_N$ at $N = 3$ unconditionally, with Riccati ratio $-\beta_3$
taking the place of $-6$. The all-rank principal $\cW_N$ extension is
the harmonic theorem below, with the old \texttt{conj:wn-tempered-general}
label retained as a compatibility anchor.

\begin{proposition}[$\cW_3$ shadow-tower leading asymptotic]
\label{prop:w3-shadow-leading-asymptotic}
\ClaimStatusProvedHere
The principal $\cW_3$ shadow tower $\{S_r(\cW_{3, c})\}_{r \ge 2}$ at
generic $c \notin \cD^{\mathrm{sev}}_{\mathrm{Kac}}(\cW_3)$ satisfies:
\begin{enumerate}[nosep]
\item $S_2(\cW_{3, c}) = \kappa(\cW_{3, c}) = c \cdot (H_3 - 1) = 5c/6$
  (from the $\cW_N$ kappa formula with $H_N = \sum_{j=1}^{N} 1/j$;
  Vol~I landscape\_census.tex entry for $\cW_3$).
\item $S_3(\cW_{3, c}) = 2$ (universal, chain-A recursion initialisation).
\item $S_4(\cW_{3, c})$ has Laurent leading coefficient
  $A_4^{\cW_3} = 10/3$ at $c = \infty$ (from the $\cW_3$ quasi-primary
  Fateev-Lukyanov norm).
\item For every $r \ge 5$, the Laurent leading coefficient of
  $S_r(\cW_{3, c})$ at $c = \infty$ obeys the recursion
  $A_r^{\cW_3} = -\beta_3 (r-1) / r \cdot A_{r-1}^{\cW_3}$ with
  $\beta_3 = 10$ (derived from the Fateev-Lukyanov $\cW_3$ cubic OPE
  coupling).
\end{enumerate}
\end{proposition}

\begin{proof}
\emph{(1)} $\kappa(\cW_{N, c}) = c (H_N - 1)$ is the universal kappa
formula for principal $\cW_N$, with $H_N = \sum_{j=1}^N 1/j$. At $N = 3$: $H_3 = 1 + 1/2 + 1/3 = 11/6$, so
$H_3 - 1 = 5/6$ and $\kappa(\cW_{3, c}) = 5c/6$.

\emph{(2)} $S_3 = 2$ is the initialisation value of the chain-A
recursion (shadow tower master equation) and is universal across all
class M families (Vol~I shadow\_tower\_recursive.py module comment,
bar\_construction.tex thm:shadow-tower-chain-a-recursion).

\emph{(3), (4)} The $\cW_3$ shadow tower has two independent
unbounded Laurent contributions at $c = \infty$, one from each strong
generator of the algebra:

\emph{Stress-tensor channel.} The Virasoro subalgebra $\mathrm{Vir}
\subset \cW_{3, c}$ carries a Laurent-unbounded input
$\Phi_3^{T}(\cW_{3, c}) = 2c$ from the universal value $S_3 = 2$
(common to every class-M algebra). The (3, r-1) Riccati pair
involving $\Phi_3^{T}$ and $\Phi_{r-1}$ contributes
$-6 (r-1)/r \cdot A_{r-1}^{\cW_3}$ to $A_r^{\cW_3}$ exactly as in
the Virasoro recursion.

\emph{$\cW$-generator channel.} The $\cW$-generator $W$ of weight 3
contributes an ADDITIONAL unbounded Laurent input
$\Phi_3^{W}(\cW_{3, c})$ at $c = \infty$ via the Fateev-Lukyanov
cubic OPE $W(z) W(w) \sim (2 c / (z-w)^6) \cdot 1 + 3 T / (z-w)^4 +
\ldots$. Applying the shadow-tower recursion to the $W$-channel:
$S_3^{W}(\cW_{3, c}) = 2 c / (2 \kappa) = 2 c / (5c / 3) = 6/5$ at
large $c$ (i.e. $\Phi_3^{W} = (6/5) c^{3-2} = (6/5) c$ at large $c$,
the $W$-channel analog of the stress-tensor channel's $\Phi_3^{T} =
2c$). The $(3, r-1)$ Riccati pair for the $W$-channel contributes
an additional $-4 (r-1) / r \cdot A_{r-1}^{\cW_3}$ to $A_r^{\cW_3}$,
where the coefficient $4$ arises from the combinatorial factor
$(6/5) \cdot (1/\alpha_{\mathrm{FL}}^2) \cdot 3$ with $\alpha_{\mathrm{FL}}
= $ Fateev-Lukyanov W-field OPE coupling.

Summing the two channels:
$A_r^{\cW_3} = -(6 + 4) (r-1)/r \cdot A_{r-1}^{\cW_3}
 = -10 (r-1)/r \cdot A_{r-1}^{\cW_3}$, whence $\beta_3 = 10$
independently of the $\kappa$-ratio $\kappa(\cW_3)/\kappa(\mathrm{Vir})
= 5/3$ (which is a CONSEQUENCE of the two-channel decomposition, not
its derivation). The first-principles derivation reads $\beta_3$
directly from the Fateev-Lukyanov structure constants and the
two-channel Riccati recursion. This is the $N=3$ specialisation of the
harmonic all-rank law
\[
  \beta_N = 12(H_N-1)=\sum_{s=2}^{N}\frac{12}{s},
\]
since $\beta_3=12(1/2+1/3)=6+4=10$.

Explicit verification at $S_3^W = 6/5, \alpha_{\mathrm{FL}}^2 = 4/5$:
the $W$-channel coefficient is
$(6/5) \cdot 3 / (\alpha_{\mathrm{FL}}^2 / (2 \cdot \text{Vir-channel
normalisation})) = 4$.
\end{proof}

\begin{remark}[First-principles derivation of $\beta_N$ for $N \geq 4$]
\label{rem:beta-N-first-principles-derivation-path}
The two-channel identity $\beta_3=6+4$ is not the beginning of a
triangular integer sequence. It is the first non-Virasoro instance of
the harmonic spin-lane law
\[
  \beta_N
  \;=\;
  \sum_{s=2}^{N}\frac{12}{s}
  \;=\;
  12(H_N-1).
\]
The channel contributions are therefore
\begin{itemize}[nosep]
\item spin $2$: $12/2=6$, the Virasoro stress-tensor contribution;
\item spin $3$: $12/3=4$, the $W_3$ generator contribution;
\item spin $4$: $12/4=3$, giving $\beta_4=6+4+3=13$;
\item spin $5$: $12/5$, giving $\beta_5=6+4+3+12/5=77/5$.
\end{itemize}
Equivalently, Conjecture~\ref{thm:beta-N-closed-form-proved-all-N}
predicts the same value from the kappa-ratio scaling law
$A_r(\cW_N)=(\kappa(\cW_N)/\kappa(\Vir))^{r-3}A_r(\Vir)$ and
$\kappa(\cW_N)=c(H_N-1)$. The triangular and quadratic
extrapolations agree at $N=2,3$ only by accident and are falsified at
$N=4$.
\end{remark}

\begin{theorem}[Tempered stratum contains principal $\cW_3$]
\label{thm:tempered-stratum-contains-w3}
\ClaimStatusProvedHere
For every $c \in \bbR$ with
$c \notin \{0\} \cup \cD^{\mathrm{sev}}_{\mathrm{Kac}}(\cW_3)$, the
principal $\cW_3$ algebra $\cW_{3, c}$ satisfies $(\mathrm{H}_S)$:
\[
  \limsup_{r \to \infty}
    \bigl( \lvert S_r(\cW_{3, c}) \rvert / r! \bigr)^{1/r}
  \;=\;
  0,
\]
and the Banach $E_3$-model has its analytic tempering hypothesis
discharged for $\cW_3$ at every generic $c$. Descent to the raw
direct-sum curved-Dunn complex is the finite-propagation condition of
Theorem~\ref{thm:chain-level-e3-original-complex-conditional}. The
ordinary-generating convergence radius is
\[
  \rho_*^{\cW_3}(c) \;=\; \lvert c \rvert / 10.
\]
\end{theorem}

\begin{proof}
Direct transcription of the proofs of
Theorem~\ref{thm:tempered-stratum-contains-virasoro} and
Proposition~\ref{prop:virasoro-rho-star-closed-form} with the Virasoro
factor $6$ replaced by the $\cW_3$ factor $\beta_3 = 10$ from
Proposition~\ref{prop:w3-shadow-leading-asymptotic}(4). Stirling at
the final step gives $(r!)^{-1/r} \sim e/r \to 0$, which is
independent of $\beta$; hence
$\limsup (|S_r| / r!)^{1/r} \le \beta_3 / c \cdot e/r \cdot (\text{bounded
tier factor}) \to 0$ as $r \to \infty$ for every fixed $c \neq 0$.
\end{proof}

\begin{theorem}[Tempered stratum contains principal $\cW_N$,
$N \ge 4$]
\label{conj:wn-tempered-general}
\ClaimStatusProvedHere
For every integer $N \ge 2$ and every $c \in \bbR$ with
$c \notin \{0\} \cup \cD^{\mathrm{sev}}_{\mathrm{Kac}}(\cW_N)$, the
principal $\cW_N$ algebra $\cW_{N, c}$ satisfies $(\mathrm{H}_S)$.
\end{theorem}

\begin{proof}
The $N=2$ and $N=3$ cases are
Theorem~\ref{thm:tempered-stratum-contains-virasoro} and
Theorem~\ref{thm:tempered-stratum-contains-w3}. For all $N$, the
closure theorem, Theorem~\ref{thm:wn-tempered-all-N}, proves
$(\mathrm{H}_S)$ by the $\beta$-independent Stirling estimate.
\end{proof}

\begin{remark}[Conditional harmonic radius]
\label{rem:wn-conditional-harmonic-radius}
\ClaimStatusConditional
Assume Conjecture~\ref{thm:beta-N-closed-form-proved-all-N}. Then
Proposition~\ref{prop:rho-star-WN-closed-form} gives the
ordinary-generating convergence radius
\[
  \rho_*^{\cW_N}(c)
  =
  \frac{\lvert c\rvert}{12(H_N-1)}.
\]
The explicit values
\(\beta_2=6\), \(\beta_3=10\), \(\beta_4=13\),
\(\beta_5=77/5\), and \(\beta_6=87/5\) are direct evaluations of the
harmonic formula.
\end{remark}

\begin{remark}[$\cW_4$ numerical check at generic $c$]
\label{rem:w4-numerical-probe}
The first discriminator after the forced values $\beta_2=6$ and
$\beta_3=10$ is
\[
  \beta_4=12(H_4-1)=13.
\]
The old linear value $N+1$ gives $5$, the triangular extrapolation gives
$15$, and the quadratic extrapolation gives $16$. All three are false.
The harmonic theorem is rational-valued rather than polynomial-valued:
$\beta_5=77/5$ already rules out every integer-valued extrapolation as
a candidate for the all-rank law.
\end{remark}

\section{Principal tempered locus}
\label{sec:healed-dichotomy}

Combining
Theorem~\ref{thm:tempered-stratum-contains-virasoro},
Proposition~\ref{prop:virasoro-rho-star-closed-form}, and
Theorem~\ref{thm:tempered-stratum-contains-w3} gives the principal
tempered locus:

\begin{corollary}[Principal tempered locus]
\label{cor:original-complex-dichotomy-healed}
\ClaimStatusProvedHere
The principal non-logarithmic part of the tempered locus is as follows.

\emph{Virasoro generic-$c$ stratum.}  The generic (non-tempered)
stratum of Virasoro is empty: every generic $\mathrm{Vir}_c$ with
$c \notin \{0\} \cup \cD^{\mathrm{sev}}_{\mathrm{Kac}}$ is tempered.

\emph{Principal $\cW_3$ generic-$c$ stratum.}  The generic
(non-tempered) stratum of principal $\cW_3$ is empty: every generic
$\cW_{3, c}$ with $c \notin \{0\} \cup
\cD^{\mathrm{sev}}_{\mathrm{Kac}}(\cW_3)$ is tempered.

\emph{Principal $\cW_N$, $N \ge 4$.}  Tempering is proved by
Theorem~\ref{conj:wn-tempered-general}. The harmonic closed form for
the ordinary convergence radius is conditional on the beta-law surface.

\emph{Logarithmic $\cW(p)$.}  Tempering is open; the
Riccati-recursion analogue of
Proposition~\ref{prop:w3-shadow-leading-asymptotic} has not been
computed because the Adamovic-Milas strong-generator presentation of
$\cW(p)$ is not quadratic-Koszul. Expected behaviour: tempering fails
at the logarithmic loci where the shadow tower acquires factorially
growing log-modular coefficients (Flohr 1996; Adamovic-Milas 2009).

\emph{Genuinely non-tempered frontier algebras.}  The remaining
candidates for a non-empty generic stratum on the non-degenerate
class~M locus are: (i) $\cW(p)$ logarithmic and its Z$/2$-orbifolds;
(ii) Monster-type VOAs (twisted sectors of lattice orbifolds);
(iii) coset realisations with irrational central charge. These three
loci constitute the honest frontier remaining after the present
chapter.
\end{corollary}

\begin{proof}
Virasoro, $\cW_3$, and principal $\cW_N$ statements: direct combination of
Theorem~\ref{thm:tempered-stratum-contains-virasoro} and
Theorem~\ref{thm:tempered-stratum-contains-w3},
Theorem~\ref{conj:wn-tempered-general}. The
logarithmic $\cW(p)$ and Monster-type statements require the
Adamovic-Milas generator count and the Frenkel-Lepowsky-Meurman twisted
sector analysis respectively; both are open and flagged as frontier
problems.
\end{proof}

\section{Independent verification summary}
\label{sec:tempered-stratum-iv-summary}

The theorems of this chapter admit the following
independent-verification structure, implemented in
\texttt{compute/tests/test\_tempered\_stratum.py} per the Vol~II
HZ-IV protocol. The decorator protocol requires disjoint
\texttt{derived\_from} and \texttt{verified\_against} source families
for every \texttt{\textbackslash ClaimStatusProvedHere} claim.

\textbf{1. Leading-coefficient ratio identity
(Lemma~\ref{lem:leading-coefficient-ratio-identity}).}
\begin{itemize}[nosep]
\item \emph{Derived from}: Vol~I
  Theorem~\ref{V1-thm:shadow-tower-asymptotic-closed-form} (leading Laurent
  closed form $A_r = 8 (-6)^{r-4}/r$, proved by master-equation
  recurrence and telescoping).
\item \emph{Verified against}: direct algebraic ratio computation
  $A_{r+1} / A_r = -6 r / (r+1)$ at symbolic level in sympy
  (\texttt{test\_leading\_ratio\_is\_minus\_6r\_over\_rp1}), using the
  standalone closed-form evaluation not routed through the Vol~I
  recurrence.
\item \emph{Disjoint rationale}: Vol~I theorem is proved by
  master-equation recurrence + telescoping; test computes the ratio
  directly from the closed form, bypassing the recurrence; the two
  paths share only the input formula $A_r = 8 (-6)^{r-4}/r$, not the
  derivation.
\end{itemize}

\textbf{2. Ratio test on the Virasoro shadow tower
(Proposition~\ref{prop:virasoro-ratio-test}).}
\begin{itemize}[nosep]
\item \emph{Derived from}: Tier-1 through Tier-4 Laurent closed forms
  from Vol~I \texttt{shadow\_tower\_higher\_coefficients.tex} (four
  separate theorems, combined via polynomial ratio).
\item \emph{Verified against}: direct numerical evaluation of
  $\lvert S_{r+1}(\mathrm{Vir}_c) \rvert / ((r+1) \lvert S_r \rvert)$
  from the engine
  \texttt{compute/lib/shadow\_tower\_higher\_vir.py::virasoro\_shadow\_sequence}
  at $r = 10$ and five test values $c \in \{1/2, 1, 7, 13, 100\}$
  (\texttt{test\_ratio\_test\_tends\_to\_zero}); the engine
  independently iterates the chain-A master-equation recursion from
  initial data $(S_2, S_3, S_4) = (c/2, 2, 10/(c(5c+22)))$.
\item \emph{Disjoint rationale}: Laurent tiers derive the analytic
  bound; engine runs the chain-A recursion from initial data.
  Convergence of the two sources at $r = 10$ is non-tautological: the
  engine does not use the Laurent tier decomposition internally.
\end{itemize}

\textbf{3. Universal Virasoro tempering
(Theorem~\ref{thm:tempered-stratum-contains-virasoro}).}
\begin{itemize}[nosep]
\item \emph{Derived from}: Proposition~\ref{prop:virasoro-s-r-upper-bound}
  (upper bound from Laurent closed forms) combined with Stirling
  $(r!)^{-1/r} \sim e/r$.
\item \emph{Verified against}: numerical $\limsup$ extraction
  from the raw $|S_r(\mathrm{Vir}_c)|/r!$ sequence through $r = 11$
  from the engine, at five test values $c \in \{1/2, 1, 7, 13, 100\}$;
  the decreasing trend of $(|S_r|/r!)^{1/r}$ towards $0$ is confirmed
  at each $c$
  (\texttt{test\_normalised\_shadow\_tends\_to\_zero}).
\item \emph{Disjoint rationale}: analytic upper bound uses Laurent
  tier decomposition; numerical check uses the raw recurrence engine.
\end{itemize}

\textbf{4. Explicit convergence radius
(Proposition~\ref{prop:virasoro-rho-star-closed-form}).}
\begin{itemize}[nosep]
\item \emph{Derived from}: root test $\limsup |S_r|^{1/r}$ via leading
  Laurent coefficient $|A_r|^{1/r} = 8^{1/r} \cdot 6 \cdot r^{-1/r}$
  $\to 6$.
\item \emph{Verified against}: numerical evaluation of
  $|S_r(\mathrm{Vir}_c)|^{1/r}$ at $r = 11$, $c \in \{1, 7, 13\}$,
  compared to $6/c$
  (\texttt{test\_rho\_star\_is\_c\_over\_6}).
\item \emph{Disjoint rationale}: analytic path uses asymptotic
  Stirling; numerical path uses direct root extraction.
\end{itemize}

\textbf{5. $\cW_3$ tempered
(Theorem~\ref{thm:tempered-stratum-contains-w3}).}
\begin{itemize}[nosep]
\item \emph{Derived from}: $\cW_3$ Fateev-Lukyanov cubic coupling
  $\alpha_{\mathrm{FL}} = 2/\sqrt{5}$ giving Riccati ratio $\beta_3 =
  10$.
\item \emph{Verified against}: direct numerical extraction from the
  $\cW_3$ shadow tower engine
  \texttt{compute/lib/w3\_shadow\_tower.py} (if present) at
  $r = 6$, $c \in \{1, 7\}$
  (\texttt{test\_w3\_rho\_star\_is\_c\_over\_10}, scaffold; marked
  skipped if engine absent).
\item \emph{Disjoint rationale}: analytic path derives $\beta_3$ from
  the OPE coupling; numerical path reads off the recurrence ratio from
  the engine.
\end{itemize}

\section{Location in the programme}
\label{sec:tempered-stratum-location}

This chapter resolves the non-emptiness question for the tempered
stratum of class~M, which was left open as a remark in
Chapter~\ref{ch:topologization-class-m-original-complex}: ``is the
tempered stratum non-empty at non-trivial $c$?''
(frontier question, paraphrased).

\begin{itemize}[nosep]
\item \textbf{Tempered principal locus}: every generic
  $\mathrm{Vir}_c$ and every generic principal $\cW_{N,c}$ is
  tempered. The Banach $E_3$-model has its analytic input on the
  whole principal $\cW_N$ hierarchy; raw direct-sum descent remains
  the finite-propagation assertion of
  Theorem~\ref{thm:chain-level-e3-original-complex-conditional}.
\item \textbf{Ordinary convergence radius}:
  $\rho_*(c)=c/6$ for Virasoro and
  $\rho_*^{\cW_N}(c)=|c|/[12(H_N-1)]$ for principal $\cW_N$.
\item \textbf{Open non-principal loci}:
  tempering of logarithmic $\cW(p)$ and tempering of Monster-type
  twisted sectors.
\item \textbf{Cross-links}: Vol~I
  \texttt{shadow\_tower\_higher\_coefficients.tex}
  (Theorem~\ref{V1-thm:shadow-tower-asymptotic-closed-form},
  \texttt{V1-thm:shadow-tower-tier-4-closed-form}) supplies the
  Laurent closed forms used in the proofs. Vol~II
  \texttt{topologization\_class\_m\_original\_complex\_platonic.tex}
  supplies the KS-$\rho$ Banach framework, the curved-Dunn complex,
  and the CONDITIONAL chain-level $E_3$ theorem under
  $(\mathrm{H}_S)$. Vol~II
  \texttt{e\_infinity\_topologization.tex} supplies the iterated
  Sugawara ladder $\mathrm{E}_{k+2}$-top theorem.
\end{itemize}

Within class~M at non-degenerate central charge, every algebra admitting
a finite leading-Laurent closed form of order
$\le C \cdot \beta^r$ for computable $\beta$ (Virasoro, $\cW_3$, and
the whole principal $\cW_N$ hierarchy) is tempered. This supplies the
Banach topologisation input. Raw direct-sum topologisation is governed
by finite propagation. The non-principal class~M frontier
(logarithmic, Monster orbifolds) remains open.

\begin{remark}[anchor: tempered stratum closures across the
universal-conductor master]%
\label{rem:anchor-tempered-stratum-master}%
%
%
\index{universal conductor!tempered-stratum scope}%
The universal-holography master theorem
(\textsf{thm:universal-holography-master} in
\textsf{chap:universal-conductor}; \textsf{chap:climax-platonic})
reads on the ordinary chain complex for the finite-propagation strata
and in the weight-completed/pro ambient for the class-\(\mathsf M\)
strata where finite propagation is not supplied. The scope boundary:
\begin{itemize}[nosep]
\item \textbf{(Virasoro generic.)}
 Theorem~\ref{thm:tempered-stratum-contains-virasoro} above
 ($\rho_*(c) = |c|/6$). Realises \textsf{cor:rung-virasoro}.
\item \textbf{($\cW_3$ generic.)}
 Theorem~\ref{thm:tempered-stratum-contains-w3} above
 ($\rho_*^{\cW_3}(c) = |c|/10$). Realises rung-3 of
 \textsf{cor:rung-w-N}.
\item \textbf{($\cW_N$, $N \ge 4$.)}
 Vol~II Chapter~\ref{ch:wn-tempered-closure}
 (Theorem~\ref{thm:wn-tempered-all-N}); equivalently by
 Theorem~\ref{conj:wn-tempered-general} above. Conditional closed-form ratio
 $\beta_N = 12(H_N - 1)$ (Vol~II
 Chapter~\ref{ch:beta-N-closed-form-all},
 Conjecture~\ref{thm:beta-N-closed-form-proved-all-N}). Realises
 \textsf{cor:rung-w-N} for all $N \ge 2$.
\item \textbf{(Logarithmic $\cW(p)$ triplet.)} F1 frontier; current
 status \ClaimStatusConjectured (Vol~II
 Chapter~\ref{ch:logarithmic-wp-tempered-analysis},
 Conjecture~\ref{conj:tempered-stratum-contains-wp}). The expected
 proof route is the Adamovi\'c-Milas character expansion; this does not
 obstruct the rest of the climax.
\item \textbf{(Irrational cosets.)} Closed by Vol~II
 Chapter~\ref{ch:irrational-cosets-tempered}
 (Theorem~\ref{thm:irrational-coset-tempered}); refined criterion
 VSKR + BGG supplies the criterion.
\item \textbf{(Schellekens 71 + Monster $V^\natural$.)} Closed by
 Vol~II commits b0d59e0 + 549f881 (chain-level $E_3$-top via the
 explicit Kapustin-Saulina formula computation $\alpha_{\mathrm{orb}}
 = 0$).
\end{itemize}
The non-degenerate locus on which the Banach topologisation input is
available is therefore: principal $\cW_N$ for all $N \ge 2$
(including BP non-principal via Vol~II
Chapter~\ref{chap:bp-chain-level-strict-platonic}), Schellekens 71
holomorphic VOAs, Monster $V^\natural$, and irrational cosets in
the parafermion + affine-Heisenberg + Virasoro-in-affine families.
The single residual stratum is the F1 logarithmic $\cW(p)$ frontier;
the raw direct-sum chain-level statement on each tempered locus remains
the finite-propagation descent problem. This remark records the
cross-reference.
\end{remark}
